% Auto-generated by reports/tables/make_data_tables.py. Do not edit.
\begin{tabular}{@{}l r r r@{}}
\toprule
 & \textbf{10-second} & 1-minute & 10-second, shorter episode \\
\midrule
bar width & 10 seconds & 1 minute & 10 seconds \\
decisions per episode & 180 & 30 & 60 \\
episode length & 30 minutes & 30 minutes & 10 minutes \\
training episodes & 22,016 & 22,276 & 66,795 \\
validation episodes & 3,885 & 3,931 & 11,787 \\
held-out test episodes & 6,474 & 6,551 & 19,645 \\
last training bar & \scriptsize 2025-07-14 03:29:50 & \scriptsize 2025-07-16 07:29:00 & \scriptsize 2025-07-15 02:59:50 \\
first test bar & \scriptsize 2025-07-14 04:00:00 & \scriptsize 2025-07-16 08:00:00 & \scriptsize 2025-07-15 03:10:00 \\
test episodes: calm / volatile & 4,559 / 1,915 & 4,596 / 1,955 & 13,636 / 6,009 \\
\midrule
\multicolumn{4}{@{}p{12.4cm}}{\footnotesize The held-out test is the last 20\% of each version's episodes and validation the last 15\% of what remains. The rule is identical across the three; because they hold different numbers of episodes, the same fraction falls on a different date in each. Between the last training bar and the first test bar sits 1 whole episode plus one bar, so no episode straddles the split. The split is written into the dataset files the agents read, so no training could have preceded it. \textbf{Episode counts quoted elsewhere in this chapter are the 10-second version's.}} \\
\bottomrule
\end{tabular}
